\documentclass[10pt,a4paper,twocolumn]{article}
\usepackage[utf8]{inputenc}
\usepackage[T1]{fontenc}
\usepackage[magyar]{babel}
\usepackage{amsmath,amssymb}
\usepackage{geometry}
\geometry{top=2cm, bottom=2cm, left=1.5cm, right=1.5cm}
\usepackage{tcolorbox}
\usepackage{xcolor}
\usepackage{booktabs}
\usepackage{multicol}

\title{deriválás, elemi függvények deriválása, aarkuszfüggvények, hiperbolikus függvények deriváltja, deriválási szabályok, inverz függvények deriválás, láncszabály}
\author{Matematika Képletgyűjtemény}
\date{}

\tcbuselibrary{skins,breakable}

\newtcolorbox{fogalombox}[1]{
  colback=blue!5!white,
  colframe=blue!75!black,
  fonttitle=\bfseries,
  title=#1,
  boxrule=0.5mm,
  arc=2mm,
  left=2mm,
  right=2mm,
  top=2mm,
  bottom=2mm,
  toptitle=1mm,
  bottomtitle=1mm
}

%%% Megjelenítéshez automatikusan kiegészített csomagok %%%
\usepackage{amsfonts, amsthm}
\usepackage{tikz}
\usepackage{pgfplots}
\pgfplotsset{compat=1.18}
\usepackage{graphicx}
\usepackage{enumitem}
\usepackage{float}
\usepackage{caption}
\usepackage{microtype}

\usepackage{hyperref}

\begin{document}

\maketitle
\vspace{-1em}
\noindent\rule{\linewidth}{0.5pt}

\section{Bevezetés és Alapfogalmak}

\begin{fogalombox}{Derivált fogalma}
A differenciálhányados (derivált) szemléletesen a függvény grafikonjának adott pontbeli érintőjének meredeksége.
\[ f'(x_0) = \frac{df}{dx}\Big|_{x=x_0} = \lim_{h \to 0} \frac{f(x_0+h) - f(x_0)}{h} \]
\textbf{Változók és dimenziók:}
\begin{itemize}
    \item $x \in \mathbb{R}$ – független változó (dimenziómentes vagy $[m]$, $[s]$ stb.)
    \item $f(x)$ – függő változó (függ a kontextustól)
    \item $h$ – kis eltolódás ($[x]$)
\end{itemize}
\end{fogalombox}

\section{Deriválási Szabályok}

\begin{fogalombox}{Alapműveletek deriváltjai}
Legyenek $f, g$ differenciálható függvények, $c \in \mathbb{R}$ konstans.
\begin{align*}
    (c \cdot f)' &= c \cdot f' \\
    (f \pm g)' &= f' \pm g' \\
    (f \cdot g)' &= f'g + fg' \quad \text{(szorzatszabály)} \\
    \left(\frac{f}{g}\right)' &= \frac{f'g - fg'}{g^2} \quad (g \neq 0, \text{ hányados})
\end{align*}
\end{fogalombox}

\begin{fogalombox}{Láncszabály és Inverz}
\textbf{Láncszabály (összetett függvény):}
\[ [f(g(x))]' = f'(g(x)) \cdot g'(x) \]
\textbf{Inverz függvény deriváltja:}
Ha $f$ invertálható és $f'(g(x)) \neq 0$:
\[ (f^{-1})'(x) = \frac{1}{f'(f^{-1}(x))} \]
\end{fogalombox}

\section{Elemi Függvények Deriváltjai}

\begin{center}
\small
\begin{tabular}{ll}
\toprule
\textbf{Függvény $f(x)$} & \textbf{Deriváltja $f'(x)$} \\
\midrule
$c$ (konstans) & $0$ \\
$x^n$ ($n \in \mathbb{R}$) & $n x^{n-1}$ \\
$e^x$ & $e^x$ \\
$a^x$ ($a > 0$) & $a^x \ln(a)$ \\
$\ln(x)$ ($x > 0$) & $\frac{1}{x}$ \\
$\log_a(x)$ & $\frac{1}{x \ln(a)}$ \\
$\sin(x)$ & $\cos(x)$ \\
$\cos(x)$ & $-\sin(x)$ \\
$\operatorname{tg}(x)$ & $\frac{1}{\cos^2(x)}$ \\
$\operatorname{ctg}(x)$ & $-\frac{1}{\sin^2(x)}$ \\
\bottomrule
\end{tabular}
\end{center}
\vspace{-1em}
*(A trigonometrikus függvényeknél a szöveg radiánban értendő).*

\section{Arkuszfüggvények Deriváltjai}

\begin{center}
\small
\begin{tabular}{ll}
\toprule
\textbf{Inverz trigl. $f(x)$} & \textbf{Deriváltja $f'(x)$} \\
\midrule
$\arcsin(x)$ ($|x|<1$) & $\frac{1}{\sqrt{1-x^2}}$ \\
$\arccos(x)$ ($|x|<1$) & $-\frac{1}{\sqrt{1-x^2}}$ \\
$\operatorname{arctg}(x)$ & $\frac{1}{1+x^2}$ \\
$\operatorname{arcctg}(x)$ & $-\frac{1}{1+x^2}$ \\
\bottomrule
\end{tabular}
\end{center}
\vspace{-1em}
\begin{fogalombox}{Arkuszfüggvények megjegyzés}
Az arkuszfüggvények az elemi trigonometrikus függvények korlátozott tartományon vett inverzei. Változók: $x$ dimenziómentes (vagy szögfüggvény értéke).
\end{fogalombox}

\section{Hiperbolikus Függvények}

\begin{center}
\small
\begin{tabular}{ll}
\toprule
\textbf{Hiperbolikus $f(x)$} & \textbf{Deriváltja $f'(x)$} \\
\midrule
$\operatorname{sh}(x)$ ($\sinh$) & $\operatorname{ch}(x)$ ($\cosh$) \\
$\operatorname{ch}(x)$ ($\cosh$) & $\operatorname{sh}(x)$ ($\sinh$) \\
$\operatorname{th}(x)$ ($\tanh$) & $\frac{1}{\operatorname{ch}^2(x)}$ \\
$\operatorname{cth}(x)$ ($\coth$) & $-\frac{1}{\operatorname{sh}^2(x)}$ \\
\bottomrule
\end{tabular}
\end{center}
\vspace{-1em}

\section{Area-függvények}

\begin{center}
\small
\begin{tabular}{ll}
\toprule
\textbf{Area-függvény $f(x)$} & \textbf{Deriváltja $f'(x)$} \\
\midrule
$\operatorname{arsh}(x)$ & $\frac{1}{\sqrt{x^2+1}}$ \\
$\operatorname{arch}(x)$ ($x>1$) & $\frac{1}{\sqrt{x^2-1}}$ \\
$\operatorname{arth}(x)$ ($|x|<1$) & $\frac{1}{1-x^2}$ \\
$\operatorname{arcth}(x)$ ($|x|>1$) & $\frac{1}{1-x^2}$ \\
\bottomrule
\end{tabular}
\end{center}

\begin{fogalombox}{Fontos azonosságok}
\begin{align*}
    \operatorname{ch}^2(x) - \operatorname{sh}^2(x) &= 1 \\
    \operatorname{th}(x) &= \frac{\operatorname{sh}(x)}{\operatorname{ch}(x)}
\end{align*}
A hiperbolikus függvények exponenciális kifejezésekkel adhatók meg:
\[ \operatorname{sh}(x) = \frac{e^x - e^{-x}}{2}, \quad \operatorname{ch}(x) = \frac{e^x + e^{-x}}{2} \]
\end{fogalombox}

\end{document}